The proceeds curve combines two opposing effects of a larger promised repayment. Holding its price fixed, an additional unit of face raises current funds. Holding resources fixed, however, it also brings more output realizations into default. Under zero recovery, investors receive nothing in those realizations. If $PD(d;c)=\Prob(d>cg)$, the associated promised gross yield can be written
\begin{equation}\label{eq:promised_yield_exposition}
 R(d;c)=\frac{R_f}{1-PD(d;c)},\qquad b(d;c)=\frac{d}{R(d;c)}.
\end{equation}
This yield compensates expected default losses under the maintained pricing assumptions. It is not an additional premium placed on top of the survival probability. The peak of the proceeds curve balances the marginal increase in the promise against the resulting loss of survival-weighted repayments. At faces beyond that peak, a larger promised payment is worth less at issuance. This is the debt Laffer-curve logic underlying the MSB definition in \textcite{CollardEtAl2015}.

Care is needed about the dates and units in this calculation. Face $d_t$ is measured relative to current GDP, whereas the surplus and next period's new borrowing are ratios to next period's GDP. Multiplication by $g_{t+1}$ puts them on the same basis as the promise. The continuation term represents the proceeds of new claims issued after next period's output is realized; it is not an additional primary surplus. Its availability depends on the future contract and pricing environment, which is why changing the instrument repeatedly affects today's resources even if the current surplus parameter is unchanged. The stationary comparison later in the paper makes this continuation channel explicit.

The distinction also explains why one-period levels should not be compared directly with observed stocks of public debt. A face value due within a five-year model period is a service obligation, while a reported public-debt stock includes claims maturing at different dates. Annualizing a cumulative surplus parameter is useful for interpreting its order of magnitude, but it does not convert this model into an annual debt-management problem. Similarly, a higher MSB measures additional proceeds at the optimal issue, while a higher MSD measures the optimizing promised repayment. Both depend on the convention used to normalize a contingent claim's face value.

The climate-growth formulation supplies the resource side of this mechanism. A transition wedge scales the growth factor when emissions change relative to their reference level. The damage wedge reflects the change in a temperature-dependent output factor between successive periods. Because these terms are multiplicative, a given climate path changes the scale of resources across the full distribution of ordinary shocks. Disaster jumps add an acute margin: two paths with similar central growth can differ in the probability of resource realizations near the default boundary. A contract comparison therefore needs the joint distribution of the payoff trigger and resources, not only a forecast of mean output.

This separation is useful when interpreting coordinated and unilateral mitigation. Global emissions enter the damage term, while a country's own transition can affect its near-term output adjustment. The baseline literature's comparison of these paths is about the evolution of fiscal resources and the value of future market access. The present paper adds a liability-design margin conditional on an assumed environment. It does not infer that changing the repayment schedule removes the real output cost of warming, or that a one-period surprise coefficient identifies an optimal mitigation path. Contract design reallocates the claims on a resource base; climate policy changes that base over time.
